% Mechanical relocation generated from the preservation ledger.
% Existing prose is included byte-for-byte from preserved blocks.

\section{Maximum Likelihood as a Reference Estimator}

ML removes the explicit prior contribution while retaining the same
parameterization and likelihood as the corresponding nonlinear MAP problem.
It is therefore the appropriate diagnostic for prior sensitivity. A large
ML--MAP difference is expected when the likelihood is weak, the prior is
informative, or the optimum lies near a bound. Conversely, numerical agreement
does not prove identification.

The public interface \texttt{compare\_reduced\_od\_likelihoods} fits Poisson
and negative-binomial specifications from the same gravity features and
response operator. The shared artifact fingerprint establishes that the
comparison changes the observation model rather than the underlying network
problem, while distinct model fingerprints prevent accidental result reuse.
Poisson is primarily a dispersion diagnostic: a material improvement from the
negative-binomial model indicates variation beyond the Poisson mean--variance
relationship, but does not by itself establish that the gravity specification
is adequate.

The companion \texttt{run\_reduced\_od\_prior\_sensitivity} operation first
computes ML and then warm-starts named MAP scenarios from the ML solution.
Each scenario records its prior fingerprint and displacement from ML. This
makes weak, moderate, and informative prior comparisons reproducible without
silently changing the likelihood or prepared response artifacts.

Because the negative-binomial raw vector contains a dispersion coordinate that
is absent from Poisson, one aligned bounds vector cannot safely configure both
fits. The comparison interface therefore accepts likelihood-specific fit
configurations. Bounds may be keyed by canonical raw parameter name and are
resolved against each actual layout before compilation. Shared parameter names
receive the same declared intervals, while the dispersion interval applies
only to negative binomial. Returned entries retain the resolved ordering,
initial point, bounds, shared prepared-artifact fingerprint, and distinct model
fingerprint. Structured callbacks expose compilation, optimization iterations,
checkpoint writes, deadlines, and model completion without adding objective
evaluations.

\subsection{Maximum likelihood estimation}

The maximum likelihood estimator (MLE) is
\[
\widehat{\bm\eta}_{\mathrm{ML}}
=
\arg\max_{\bm\eta}
\ell(\bm\eta),
\]
where
\[
\ell(\bm\eta)
=
\log\mathcal L\!\left(
Y\mid \bm f(\widetilde{\bm z}),\theta(\widetilde u),\rho,r
\right),
\]
with the dependence on $\widetilde u$ omitted when $\theta$ is fixed. The MLE
uses no prior information. The baseline demand $\bm f^{(0)}$ still appears in
the parameterization $\bm f=\bm f^{(0)}\exp(\bm z)$, but it is not a probabilistic
prior in the ML objective.

After optimization, the raw point estimate is transformed using exactly the same
smooth mappings as in Bayesian inference:
\[
\widehat{z}_q^t
=z_{\max}\tanh(\widehat{\widetilde z}_q^t/z_{\max}),
\qquad
\widehat f_q^t=f_q^{t,(0)}\exp(\widehat z_q^t),
\]
and, when estimated,
\[
\widehat u
=u_{\max}\tanh(\widehat{\widetilde u}/u_{\max}),
\qquad
\widehat\theta=\exp(\widehat u).
\]


